
\DocumentMetadata{lang=en-US,testphase=phase-III}
\documentclass[10pt]{article}
%
% LaTeX Problem Set Template originally designed by Stanford CS103.

\usepackage{amsmath}
\usepackage{amssymb}
\usepackage{amsthm}
\usepackage{amssymb}
\usepackage{mathdots}
\usepackage{braket}
\usepackage[pdftex]{graphicx}
\usepackage{fancyhdr}
\usepackage[margin=1in]{geometry}
\usepackage{multicol}
\usepackage{bm}
\usepackage{listings}
\PassOptionsToPackage{usenames,dvipsnames}{color}  %% Allow color names
\usepackage{pdfpages}
\usepackage{algpseudocode}
\usepackage{tikz}
\usepackage{tikzsymbols}                           %% jj addition
\usepackage{enumitem}
\usepackage[T1]{fontenc}
\usepackage{inconsolata}
\usepackage{framed}
\usepackage{wasysym}
\usepackage[thinlines]{easytable}
\usepackage{hyperref}
\usepackage{wrapfig}
\hypersetup{
    colorlinks=true,
    linkcolor=blue,
    filecolor=magenta,
    urlcolor=blue,
    title={HW1 CS380 Fall 2026}
    pdftitle={HW1 CS380 Fall 2026},
    pdfauthor={J.K. Johnstone},
    pdfsubject={Vector},
    pdfkeywords={vector, inner product}
}

\title{CS 380: Math for AI; Homework \#1: The vector}
\author{[TODO: Replace this with your full name]}
\date{\today}

% Running author based on https://tex.stackexchange.com/questions/68308/how-to-add-running-title-and-author#answer-68310
\makeatletter
\let\runauthor\@author
\makeatother

\lhead{\runauthor}
\chead{Homework 1}
\rhead{\today}
\lfoot{}
\cfoot{CS 380: Matrix Computation --- Fall 2025}
\rfoot{\thepage}

\newcommand{\abs}[1]{\lvert #1 \rvert}
\renewcommand{\(}{\left(}
\renewcommand{\)}{\right)}
\newcommand{\floor}[1]{\left\lfloor#1\right\rfloor}
\newcommand{\ceil}[1]{\left\lceil#1\right\rceil}
\newcommand{\pd}[1]{\frac{\partial}{\partial #1}}
\newcommand{\powerset}[1]{\wp\left(#1\right)}
\newcommand{\suchthat}{\ \vert \ }
\newcommand{\naturals}{\mathbb{N}}
\newcommand{\integers}{\mathbb{Z}}
\newcommand{\reals}{\mathbb{R}}
\renewcommand{\qed}{\blacksquare}
\newcommand{\accepts}{\text{ accepts }}
\newcommand{\rejects}{\text{ rejects }}
\newcommand{\loopson}{\text{ loops on }}
\newcommand{\haltson}{\text{ halts on }}
\newcommand{\encoded}[1]{\left\langle#1\right\rangle}
\newcommand{\rlangs}{\mathbf{R}}
\newcommand{\relangs}{\mathbf{RE}}
\newcommand{\corelangs}{\text{co-}\mathbf{RE}}
\newcommand{\plangs}{\mathbf{P}}
\newcommand{\nplangs}{\mathbf{NP}}
\renewcommand{\labelitemii}{$\bullet$}
\renewcommand\qedsymbol{$\blacksquare$}
\newenvironment{prf}{{\bfseries Proof.}}{\qedsymbol}
\renewcommand{\emph}[1]{\textit{\textbf{#1}}}
\newcommand{\annotate}[1]{\textit{\textcolor{blue}{#1}}}
\usepackage{mdframed}
\usepackage{float}

\makeatother

\definecolor{shadecolor}{gray}{0.95}

\theoremstyle{plain}
\newtheorem*{lem}{Lemma}
\newtheorem{defn2}{Definition}                              % jj addition

\theoremstyle{plain}
\newtheorem*{claim}{Claim}

\theoremstyle{definition}
\newtheorem*{answer}{Answer}

\newtheorem{theorem}{Theorem}[section]
\newtheorem*{thm}{Theorem}
\newtheorem{corollary}{Corollary}[theorem]
\newtheorem{lemma}[theorem]{Lemma}

\renewcommand{\headrulewidth}{0.4pt}
\renewcommand{\footrulewidth}{0.4pt}

\setlength{\parindent}{0pt}

\pagestyle{fancy}

\renewcommand{\thefootnote}{\fnsymbol{footnote}}

\usepackage{boxedminipage}
\newenvironment{blank}{\colorbox{shadecolor}{\strut \underline{\ \ \ \ \ }}}

% ==================== IGNORE LATEX ABOVE EXCEPT FOR TODO ======================

\begin{document}

\maketitle

\begin{center}
  \emph{Due Monday, September 14 at 11:59pm Central}
\end{center}

% \vspace{1cm}

Please read the following independent completion attestation before starting the
homework.
Then digitally sign it,
acknowledging that you understand the guidelines for this homework,
and have followed them.

\begin{verbatim}
I declare that I have completed this assignment entirely on my own,
and following the generative AI policies.
I have read the UAB Academic Integrity Code and understand that any breach 
of this code may result in severe penalties, including failure of the class.
\end{verbatim}

Signature: 

\section*{Symbols Reference}
Here are some symbols that may be useful for this problem set\\
(see the corresponding \LaTeX\ in the raw document):

\begin{itemize}
\item summation and inner product: $a \cdot b = \sum_{i=1}^n a_i * b_i$
\item square root: $\sqrt{2}$
\item norm: $\|v\|_2$
\item fraction: $\frac{1}{2}$
\item $\theta$, as in $\cos \theta = a \cdot b$ (when both unit)
\item vector:
\[
\begin{pmatrix}
0 & 1 & 2
\end{pmatrix}
\]

\item matrix:
\[
\begin{pmatrix}
0 & 1 & 2\\
3 & 4 & 5
\end{pmatrix}
\]
\item for pretty printing, you may want to add a newline using backslash backslash (see Q1 for an example)
\item to force a new page to be begun (especially if LaTeX is acting strangely at the end of a page), use the newpage command (see a few lines below)
\end{itemize}

\LaTeX\ typing tips:
\begin{itemize}
    \item Exponents (use curly braces if exponent is more than 1 character): 
          $x^2$, $2^{3x}$
    \item Subscripts (use curly braces if subscript is more than 1 character): 
          $x_0$, $x_{10}$
    \item {\bf Boldface} and {\em italics}
\end{itemize}

{\bf All written answers will be added to this TeX document.}\\
Replace all 'Write your answer here $\ldots$' by your answer,
removing the 'Write your answer here $\ldots$'.

\newpage

% ==================== IGNORE LATEX ABOVE EXCEPT FOR TODO ======================
% ========================= START HERE =========================================

\section*{Q1: The hyperplane (written)}

    In this question, you will prove one of the key steps 
    in moving from the algebra of a hyperplane to the geometry of a hyperplane:
    the coefficient vector of the linear equation is the normal.

    \begin{shaded}
      Lemma: Let $a$ be a vector in d-space, and $b$ a scalar.\\
      The normal of the hyperplane defined by the linear
      equation $a \cdot x = b$ is $a$.\\
      \begin{prf} \\
      {\em Write your proof here.}\\
      \end{prf}
    \end{shaded}

    An oriented hyperplane is a hyperplane where one halfspace is positive
    and the other halfspace is negative, with the normal pointing towards
    the positive halfspace.
    So flipping the normal yields a different oriented hyperplane.
    One geometric representation for an oriented hyperplane
    is by a unit normal (pointing in the direction of the positive
    hyperplane) and a point $P$ on the hyperplane.

    \begin{itemize}
    \item Find a problem with this representation.
    \item Explain the problem.
    \item Find a geometric representation 
          of an oriented hyperplane that solves this problem,
          (hint: keep the normal, change the point to something else).
    \item Find another advantage of this new representation 
          (aside from solving the old problem).
    \end{itemize}

    \begin{shaded}
     {\em Write your answer here.}
    \end{shaded}

\clearpage

\section*{Q2: The word vector (written and code)}

    To prepare for this question, download the 
      glove.2024.wikigiga.50d dataset from the Stanford GloVe
      website (link in lecture), extract it, and rename it 
      'glove.2024.wikigiga.50d.txt'.
      You will use this in your experiments.\\

      For simplicity, you will use naive algorithms to work with the vectors 
      in this homework,
      since our focus is on learning how basic linear algebra of vectors.
      You can contemplate more sophisticated algorithms for future.

\section*{Q2a: Extraction of the word vectors}

      You will start by extracting the following word vectors 
      from glove.2024.wikigiga.50d.txt:
\[
      \text{good, evil, devil, angel, goods, cup, happy, ethical}
\]

      Store in a matrix, one word vector per column.
      (We call this a column-major data matrix.)\\
      The shape of the matrix will be (50,8).\\
      Use the order of the words, so the 4th column will be the 
      word vector for 'angel'.\\

      To solve this problem,
      write a Python function in your script
      that extracts a word into a NumPy vector:
      input is the path to the dataset (a string) and the word (a string);\\
      output is the word vector (NumPy array).\\
      A typical call would be:
      extract ('glove.2024.wikigiga.50d.txt', 'good')
      Use this code to build the matrix and report it below.

      Note: until we cover matrices, you could temporarily store 
      in 8 separate NumPy vectors.

      \begin{shaded}
      {\em Paste the data matrix here.}\\
       Please prettyprint the matrix: 
       one row per line, scalars nicely separated, columns lining up.\\
       The Python script should call the extraction function 
       to build this matrix.
      \end{shaded}

\section*{Q2b: Cosine similarity}

      Next find the cosine similarity 
      between each pair of word vectors,
      and store in a matrix of shape (8,8).\\
      Report each similarity to 5 decimal places.\\
      Use the above order: for example, the similarity between
      'devil' and 'angel' would be stored in the
      entry in the 3rd row and 4th column, 
      and the entry in the 4th row and 3rd column
      (hopefully they are the same!).\\

      To solve this problem,
      write a Python function in your script
      that computes cosine similarity:\\
      input is two word vectors (NumPy array, NumPy array);\\
      output is the cosine similarity (NumPy scalar).\\
      Write two versions of this function: using NumPy 
      and computing it yourself using a for loop (see the template script).
      
      \begin{shaded}
      {\em Paste the matrix output here.}\\
      Please prettyprint the matrix.\\
      The Python script should call the cosine similarity function 
      to build this matrix.
      \end{shaded}

      Which two words are most similar? least similar?
      \begin{shaded}
      Most similar word pair:  \\
      Least similar word pair: 
      \end{shaded}

\clearpage

\section*{Q2c: Euclidean length}

      Now find the length of our 8 word vectors,
      and store the results in a vector (again using the given order).\\
      This output vector will be a vector in 8-space.\\
      Report each length to 5 decimal places.\\

      To solve this problem,
      write a Python function in your script
      that computes the length of a vector in d-space:\\
      input is a vector in d-space (NumPy array); \\
      output is its Euclidean (2-norm) length (NumPy scalar).\\
      Write two versions of this function: 
      using NumPy and computing it yourself using a for loop.

      \begin{shaded}
      {\em Paste the length vector here.}
      \end{shaded}

\section*{Q2d: Mikolov example}

      One of the examples that (Mikolov 2013, 
      'Distributed Representations of Words and Phrases and their Compositionality' on arXiv)
      uses to illustrate that their word vectors behave properly
      (like vectors in linear algebra) is:
      vec('King') - vec('Man') + vec('Woman') is very close to vec('Queen').\\

      In the GloVe 50-d dataset,
      what is the cosine similarity and Euclidean distance
      between these two vectors?\\
      Report to 5 decimal places.\\
      Note that you would not expect as good results 
      with these 50-dimensional vectors of GloVe
      as the 640-dimensional vectors of Mikolov.

      \begin{shaded}
      {\em Report the similarity score and distance here.}\\
      Cosine similarity:  \\
      Distance:
      \end{shaded}

\section*{What to hand in}

Generate a pdf of this document using Overleaf,
naming the file 'hw1-LASTNAME-26fa380.pdf'.\\
Example: hw1-johnstone-26fa380.pdf\\[1em]
Write your code in a Python script named 'hw1-LASTNAME-26fa380.py'.\\
Example: hw1-johnstone-26fa380.py \\
I have given you a starting point in template-hw1-26fa380.py.\\
Please use this code as a starting point (e.g., use the same function names).\\[1em]
Zip the pdf and Python script into hw1-LASTNAME-26fa380.zip, 
and upload this one zip file to Canvas.

\end{document}
